
% !TEX program = xelatex
\documentclass[UTF8,12pt,a4paper]{ctexart}
\usepackage{geometry}
\usepackage{array,booktabs,longtable,multirow}
\usepackage{amsmath,amssymb,amsthm}
\usepackage{enumitem}
\usepackage{xcolor}
\usepackage{fancyhdr}
\usepackage{lastpage}
\usepackage{eso-pic}
\usepackage{tikz}
\geometry{left=2cm,right=2cm,top=2cm,bottom=1.8cm}
\setmainfont{Noto Serif CJK SC}
\setCJKmainfont{Noto Serif CJK SC}
\pagestyle{fancy}
\fancyhf{}
\fancyhead[C]{\small 【课程名称】期末模拟卷统一模板}
\fancyfoot[C]{\small 第 \thepage\ 页 / 共 \pageref{LastPage} 页}
\renewcommand{\headrulewidth}{0pt}
\renewcommand{\footrulewidth}{0pt}
\AddToShipoutPictureBG{%
  \begin{tikzpicture}[remember picture,overlay]
    \node[rotate=35,scale=1.6,text opacity=0.10] at (current page.center)
    {Jerry 制作｜仅供个人复习使用｜禁止盗印与商用传播};
  \end{tikzpicture}%
}
\setlist[enumerate,1]{leftmargin=2em,itemsep=0.4em}
\newcommand{\examtitle}[2]{%
\begin{center}
  {\LARGE\bfseries 《#1》本科模拟卷#2}\\[0.4em]
  {\small {{学校/学院}}\quad {{学年学期}}\quad {{专业/年级}}}\\[0.8em]
\end{center}}
\newcommand{\examinfo}{%
\begin{center}
\begin{tabular}{|>{\centering\arraybackslash}p{2.5cm}|>{\centering\arraybackslash}p{3.5cm}|>{\centering\arraybackslash}p{2.5cm}|>{\centering\arraybackslash}p{3.5cm}|}
\hline
考试方式 & {{闭卷/开卷}} & 考试时间 & {{120 分钟}} \\\hline
卷面总分 & {{100 分}} & 适用层级 & {{A/B/其他}} \\\hline
\end{tabular}
\end{center}}
\newcommand{\scoretable}{%
\begin{center}
\begin{tabular}{|c|c|c|c|c|c|}
\hline
题号 & 一 & 二 & 三 & 四 & 总分 \\\hline
得分 & & & & & \\\hline
\end{tabular}
\end{center}}
\newenvironment{choices}{\begin{enumerate}[label=(\Alph*),leftmargin=3.5em,itemsep=0.1em]}{\end{enumerate}}
\begin{document}
\examtitle{{课程名称}}{{一/二}}
\examinfo
\scoretable

\section*{一、单项选择题（每小题 3 分，共 30 分）}
\begin{enumerate}
\item 【题干】（\quad）
\begin{choices}\item 【选项A】\item 【选项B】\item 【选项C】\item 【选项D】\end{choices}
\item 【题干】（\quad）
\begin{choices}\item 【选项A】\item 【选项B】\item 【选项C】\item 【选项D】\end{choices}
\item 【题干】（\quad）
\begin{choices}\item 【选项A】\item 【选项B】\item 【选项C】\item 【选项D】\end{choices}
\item 【题干】（\quad）
\begin{choices}\item 【选项A】\item 【选项B】\item 【选项C】\item 【选项D】\end{choices}
\item 【题干】（\quad）
\begin{choices}\item 【选项A】\item 【选项B】\item 【选项C】\item 【选项D】\end{choices}
\item 【题干】（\quad）
\begin{choices}\item 【选项A】\item 【选项B】\item 【选项C】\item 【选项D】\end{choices}
\item 【题干】（\quad）
\begin{choices}\item 【选项A】\item 【选项B】\item 【选项C】\item 【选项D】\end{choices}
\item 【题干】（\quad）
\begin{choices}\item 【选项A】\item 【选项B】\item 【选项C】\item 【选项D】\end{choices}
\item 【题干】（\quad）
\begin{choices}\item 【选项A】\item 【选项B】\item 【选项C】\item 【选项D】\end{choices}
\item 【题干】（\quad）
\begin{choices}\item 【选项A】\item 【选项B】\item 【选项C】\item 【选项D】\end{choices}
\end{enumerate}

\section*{二、填空题（每小题 3 分，共 15 分）}
\begin{enumerate}
\item 【题干】答案：\underline{\hspace{3cm}}。
\item 【题干】答案：\underline{\hspace{3cm}}。
\item 【题干】答案：\underline{\hspace{3cm}}。
\item 【题干】答案：\underline{\hspace{3cm}}。
\item 【题干】答案：\underline{\hspace{3cm}}。
\end{enumerate}

\section*{三、计算题（每小题 11 分，共 44 分）}
\begin{enumerate}
\item 【题干】
\begin{enumerate}[label=(\arabic*)]
\item 【小问】
\item 【小问】
\item 【小问】
\end{enumerate}
\item 【题干】
\begin{enumerate}[label=(\arabic*)]
\item 【小问】
\item 【小问】
\item 【小问】
\end{enumerate}
\item 【题干】
\begin{enumerate}[label=(\arabic*)]
\item 【小问】
\item 【小问】
\item 【小问】
\end{enumerate}
\item 【题干】
\begin{enumerate}[label=(\arabic*)]
\item 【小问】
\item 【小问】
\item 【小问】
\end{enumerate}
\end{enumerate}

\section*{四、证明题（共 11 分）}
\begin{enumerate}
\item 【证明题干】
\end{enumerate}

\newpage
\section*{答案与解析}
\subsection*{一、单项选择题}
\begin{enumerate}
\item 答案：【A/B/C/D】。解析：【判断依据】。
\item 答案：【A/B/C/D】。解析：【判断依据】。
\item 答案：【A/B/C/D】。解析：【判断依据】。
\item 答案：【A/B/C/D】。解析：【判断依据】。
\item 答案：【A/B/C/D】。解析：【判断依据】。
\item 答案：【A/B/C/D】。解析：【判断依据】。
\item 答案：【A/B/C/D】。解析：【判断依据】。
\item 答案：【A/B/C/D】。解析：【判断依据】。
\item 答案：【A/B/C/D】。解析：【判断依据】。
\item 答案：【A/B/C/D】。解析：【判断依据】。
\end{enumerate}
\subsection*{二、填空题}
\begin{enumerate}
\item 答案：【结果】。解析：【关键计算过程】。
\item 答案：【结果】。解析：【关键计算过程】。
\item 答案：【结果】。解析：【关键计算过程】。
\item 答案：【结果】。解析：【关键计算过程】。
\item 答案：【结果】。解析：【关键计算过程】。
\end{enumerate}
\subsection*{三、计算题}
\begin{enumerate}
\item 解析：【完整推导】。
\item 解析：【完整推导】。
\item 解析：【完整推导】。
\item 解析：【完整推导】。
\end{enumerate}
\subsection*{四、证明题}
\begin{enumerate}
\item 证明：【完整证明过程】。
\end{enumerate}
\end{document}
